\section{Las cuentas de ingresos de OpenAI: ingresos visibles e ingresos potenciales no visibles}

La sección anterior explicó por qué no es fácil que los gigantes se levanten de la mesa de juego; esta sección vuelve a enfocar a OpenAI. La mayor controversia en torno a OpenAI es si sus enormes compromisos de capacidad de cómputo corresponden a su flujo de caja futuro. La forma en que Guangmi (广密) desglosa el tema resulta muy didáctica: primero separa los «ingresos visibles» y después los «ingresos no visibles». Los primeros se parecen a los de las plataformas de internet y los servicios en la nube; los segundos dependen de que los Agent puedan sustituir trabajo humano y de que la IA pueda realizar tareas de alto valor que las personas no pueden hacer.

\begin{figure}[H]
\centering
\includegraphics[width=0.92\textwidth]{openai-revenue-stack.png}
\caption{Pila de ingresos de OpenAI: suscripciones, publicidad y comercio electrónico, API, trabajo de los Agent y científicos de IA. Diagrama conceptual de elaboración propia, basado en la conversación de 00:07:38--00:13:10 y 00:13:32--00:15:04.}
\end{figure}

\begin{knowledgebox}{Lectura de la figura: de los «ingresos de plataforma de internet» al «presupuesto de trabajo»}
Las tres primeras capas son ingresos relativamente visibles: suscripciones, publicidad y comercio electrónico, y API. Las dos últimas son más difíciles de estimar: si los Agent pueden sustituir de principio a fin el trabajo de oficina, captan el presupuesto de salarios; si la IA puede escribir CUDA, hacer descubrimientos científicos o buscar nuevos materiales, crea un valor incremental que hoy las personas no han logrado producir plenamente.
\end{knowledgebox}

\subsection{Ingresos visibles: suscripciones, publicidad y comercio electrónico, y API}

La primera categoría de ingresos de Guangmi son las suscripciones. Puede expresarse con una fórmula simplificada:
\[
R_{\text{sub}} = U \times p \times a
\]
donde \(R_{\text{sub}}\) representa el ingreso anual por suscripciones; \(U\), el número de usuarios activos mensuales o de usuarios alcanzables; \(p\), la tasa de conversión a pago; y \(a\), el ingreso anual promedio por usuario de pago. La estimación optimista que se da en el programa es la siguiente: si ChatGPT llega a miles de millones de usuarios activos y alrededor del 10\% paga, los ingresos por suscripciones pueden alcanzar el orden de decenas de miles de millones de dólares. Esta analogía proviene de productos de suscripción consolidados como Office y Netflix.

La segunda categoría son la publicidad y el comercio electrónico. La lógica es que un punto de acceso con mucho tráfico, uso frecuente y alta retención termina por monetizarse. En sus inicios, los videos cortos y la publicidad en teléfonos móviles también se cuestionaron porque la pantalla era pequeña y el contexto no parecía adecuado, pero después se demostró que la publicidad y las transacciones se trasladan adonde está la atención de los usuarios. Si ChatGPT integra búsqueda, recomendaciones, compras e intermediación de servicios, podría quedarse con una parte de la publicidad de búsqueda tradicional y de las comisiones de las plataformas.

La tercera categoría es la API, es decir, la «nube de modelos». Guangmi la compara con el surgimiento de AWS a partir del comercio electrónico de Amazon: cuando una empresa necesita para sí misma enormes sistemas de cómputo, datos e ingeniería, esa misma capacidad puede ofrecerse como servicio a los desarrolladores. Para OpenAI, cuanto más capaz sea el modelo, menor su costo y mayor su ecosistema de desarrolladores, más se parecerán los ingresos de la API a un nuevo negocio en la nube.

\begin{importantbox}{La condición común de los ingresos visibles}
Las suscripciones, la publicidad y el comercio electrónico, y la API dependen de lo mismo: ChatGPT debe mantener mucho tráfico, uso frecuente, una marca fuerte y una experiencia de modelo suficientemente buena. Si el punto de acceso de los usuarios se traslada a Gemini, a los fabricantes de teléfonos o a Agent verticales, el techo de estas tres categorías de ingresos bajará.
\end{importantbox}

\subsection{Ingresos no visibles: Agent que sustituyen trabajo humano}

La expectativa más grande proviene del presupuesto de trabajo. El software de herramientas vende eficiencia; si un Agent realmente completa tareas de principio a fin, lo que vende podría ser «la sustitución de una parte de los salarios». El desarrollo de software es el primer sector en el que aparece esta señal: AI Coding generó en un año un mercado del orden de diez mil millones de dólares, lo que muestra que en el trabajo del conocimiento sí existen tareas que los modelos pueden absorber con rapidez.

Pero aquí hay dos umbrales decisivos. Primero, el Agent debe ser lo bastante confiable para asumir responsabilidades, no solo para generar sugerencias; segundo, las empresas deben estar dispuestas a trasladar presupuesto de personal, subcontratación o software tradicional a sistemas de IA. Solo al superar estos dos umbrales la IA deja de limitarse a «ayudar a los empleados a trabajar más rápido» y se convierte en una nueva oferta de trabajo.

\begin{warningbox}{El presupuesto para herramientas y el presupuesto de salarios no son lo mismo}
Que un empleado esté dispuesto a pagar unas decenas de dólares al mes por una herramienta no significa que la empresa esté dispuesta a entregarle a un Agent el presupuesto de un puesto. Lo segundo exige estabilidad, control de permisos, auditoría, asignación de responsabilidades y resultados predecibles. Muchas empresas de IA se quedarán estancadas en ser «una herramienta muy útil» en lugar de «un sistema capaz de captar el presupuesto de salarios».
\end{warningbox}

\subsection{Ingresos más lejanos: la IA hace lo que las personas no pueden}

Guangmi también menciona otra categoría de ingresos no visibles: la IA podría completar tareas de alto valor que hoy las personas no pueden hacer o hacen muy lentamente, por ejemplo, reescribir CUDA, descubrir nuevos materiales, acelerar la investigación médica, buscar pistas sobre la superconductividad a temperatura ambiente e incluso «que la IA sea científica de IA». Este tipo de ingresos no puede clasificarse simplemente como sustitución de lo existente, porque podría crear nuevos problemas, nuevos mercados y nuevas rutas científicas.

Aquí la condición es más exigente: el modelo no solo debe saber utilizar el conocimiento existente, sino también mejorar mediante experimentación a largo plazo, retroalimentación, razonamiento, memoria y aprendizaje autónomo. Esto conduce de forma natural al Online Learning que se trata más adelante. Dicho de otro modo, cuanto mayor sea la expectativa de ingresos a largo plazo de OpenAI, más altas serán las exigencias de confiabilidad del modelo, aprendizaje continuo y retroalimentación de entornos reales.

\subsection{Resumen de la sección}

Las cuentas de ingresos de OpenAI pueden dividirse en tres niveles: las suscripciones, la publicidad y el comercio electrónico, y la API son ingresos de plataforma visibles; la sustitución del trabajo de oficina por Agent es una migración más grande de valor existente; y los descubrimientos científicos de la IA y las tareas incrementales de alto valor dependen de avances más fundamentales en la capacidad de los modelos. Las dudas a corto plazo provienen de las cuentas financieras; el optimismo a largo plazo, de las cuentas de plataforma y de paradigma.
